% Generated draft: requires mathematical and source review.
\documentclass[11pt]{article}
\usepackage[margin=25mm]{geometry}
\usepackage{fontspec}
\setmainfont{Latin Modern Roman}
\usepackage{amsmath,amssymb,mathtools,unicode-math}
\setmathfont{Latin Modern Math}
\usepackage{enumitem,needspace,xurl,hyperref}
\setlength{\parindent}{0pt}
\setlength{\parskip}{5pt}
\setlength{\emergencystretch}{3em}
\newcommand{\N}{\mathbb{N}}
\newcommand{\Z}{\mathbb{Z}}
\newcommand{\Q}{\mathbb{Q}}
\newcommand{\R}{\mathbb{R}}
\newcommand{\C}{\mathbb{C}}
\newcommand{\F}{\mathbb{F}}
\newcommand{\A}{\mathbb{A}}
\newcommand{\PP}{\mathbb{P}}
\newcommand{\E}{\mathbb{E}}
\newcommand{\eps}{\varepsilon}
\newcommand{\dd}{\,\mathrm{d}}
\begin{document}
\section*{UnsolvedMath ID: 700985}
\textbf{Source problem code:} catalogue entry
\subsection{Definitions and mathematical statement}
An involution in a group \(G\) is an element \(i\) of order exactly \(2\). Write \(i^g=g^{-1}ig\) and \(C_G(i)=\{x\in G:xi=ix\}\). A group is periodic if every element has finite order. A group \(H\) is residually periodic if, for every \(h\in H\setminus\{1\}\), there exists a normal subgroup \(N\trianglelefteq H\) such that \(h\notin N\) and \(H/N\) is periodic; explicitly,
\[
\forall x\in H\;\exists n\in\Z_{\geq 1}\quad (xN)^n=N.
\]
The supplied source refers to the definition of a \(T_0\)-group in Kourovka Notebook Problem 12.101, but does not supply that definition. Thus the hypothesis that \(G\) is a \(T_0\)-group remains unresolved here, and the question below is conditional on the intended meaning of that term.

For every \(T_0\)-group \(G\) and every involution \(i\in G\) such that
\[
G=\left\langle i^g\mid g\in G\right\rangle,
\]
must \(C_G(i)\) be residually periodic? Equivalently, under these hypotheses, must every \(h\in C_G(i)\setminus\{1\}\) survive in some periodic quotient of \(C_G(i)\)?
\subsection{Short English statement}
V. P. Shunkov asks whether the centralizer of an involution whose conjugates generate a \(T_0\)-group is residually periodic. The source's definition of a \(T_0\)-group is missing, so the intended class of groups requires clarification from Problem 12.101.
\subsection{Sources}
\begin{itemize}
\item Local UnsolvedMath registry, permanent ID \(700985\), exact source code \texttt{catalogue entry}, title ``Kourovka Notebook Problem 13.67''; supplied upstream revision \texttt{d371ee917ac9647fbb06cd6fc529deec6f3dda82}.
\item Supplied attribution: V. P. Shunkov, proposed in 1995; \emph{The Kourovka Notebook (Unsolved Problems in Group Theory)}, 13th issue (1995), Problem 13.67. The supplied record identifies the 21st issue edition (2026), \texttt{arXiv:1401.0300v48}. This bibliographic attribution is supplied source data and has not been independently verified.
\item The statement refers to Kourovka Notebook Problem 12.101 for the definition of a \(T_0\)-group; that referenced content is absent from the supplied record.
\item The registry labels this problem ``open''. This is a supplied status label, not an independently verified assessment of the literature.
\end{itemize}
\end{document}
